\section{Operadores de destrucción}
\label{sec:destruccion}

Los seis operadores tienen la firma $\textit{des}(x,q)$: extraen aproximadamente $q$
clientes de las rutas de $x$ y los añaden a $\Uns(x)$, que en la fase~1 puede contener ya
clientes pendientes de iteraciones anteriores (el banco de la Sección~\ref{sec:fases}). Ninguno
elimina rutas: una ruta que se queda sin clientes permanece en la solución como ruta vacía.
Todos terminan actualizando $K(x)$ y $D(x)$. Como salvaguarda, los bucles de las extracciones
aleatoria, de rutas, de peores y de Shaw se detienen si $|\Uns|$ alcanza $n-1$.

Se denomina \emph{ruta activa} a la que tiene al menos un cliente, y \emph{costo de retiro} de
un cliente $\pi_k$ al ahorro de distancia que produce quitarlo de su ruta:
\begin{equation}
\operatorname{cr}(\pi_k)=d_{\pi_{k-1}\pi_k}+d_{\pi_k\pi_{k+1}}-d_{\pi_{k-1}\pi_{k+1}}\ \ge 0 .
\label{eq:cr}
\end{equation}

\subsection{Extracción aleatoria}

\begin{algorithm}[H]
\caption{\fn{ExtracciónAleatoria} (\cod{randomRemoval})}
\label{alg:random}
\Entrada{solución $x$, grado $q$}
\Mientras{se hayan extraído menos de $q$ clientes \Y existan rutas activas}{
  elegir una ruta activa $\pi$ uniformemente al azar\;
  elegir una posición $k\sim\unifint{1}{L-2}$ de $\pi$\;
  mover $\pi_k$ a $\Uns$;\quad \fn{RecalcularRuta}($\pi$)\;
}
\end{algorithm}

El sorteo es en dos etapas (primero la ruta, luego el cliente), por lo que no es uniforme
sobre los clientes: un cliente de una ruta con $L-2$ clientes se elige con probabilidad
$1/(m_{\mathrm{act}}(L-2))$, donde $m_{\mathrm{act}}$ es el número de rutas activas, y los
clientes de rutas cortas son más propensos a salir. Es el operador de diversificación pura y
su costo es
$\bigO\bigl(q\,(m+\bar{L})\bigr)$.

\subsection{Extracción de rutas}

\begin{algorithm}[H]
\caption{\fn{ExtracciónDeRutas} (\cod{routeRemoval})}
\label{alg:route}
\Entrada{solución $x$, grado $q$}
\Mientras{se hayan extraído menos de $q$ clientes \Y existan rutas activas}{
  elegir una ruta activa $\pi$ uniformemente al azar\;
  mover todos los clientes de $\pi$ a $\Uns$;\quad $\pi\gets(0,0)$\;
}
\end{algorithm}

Vacía rutas enteras hasta alcanzar $q$, por lo que puede superar $q$ en hasta una ruta. La
ruta vaciada permanece en la solución como ruta vacía: no cuenta como vehículo, pero queda
disponible para la reparación. El operador obliga a redistribuir por completo los clientes de
varios vehículos. Como la reparación no abre rutas, es además el principal mecanismo por el
que la fase~2 puede terminar con menos vehículos que los heredados de la fase~1: basta con
que los clientes de una ruta vaciada se recoloquen en las demás.

\subsection{Extracción de los peores}
\label{sec:worst}

Extrae clientes con costo de retiro alto, es decir, los que están «mal colocados». La
elección se aleatoriza como en~\cite{ropke2006}: con la lista $\mathcal{L}$ de clientes
ordenada por costo de retiro decreciente, se toma el elemento de índice
\begin{equation}
\iota=\bigl\lfloor y^{\,p}\,|\mathcal{L}|\bigr\rfloor,\qquad y\sim\unif{0}{1},\quad p=3 .
\label{eq:sesgo}
\end{equation}
Como $P(\iota<k)=(k/|\mathcal{L}|)^{1/p}$, con $p=3$ el 10\,\% superior de la lista se elige
con probabilidad $0.46$ y el 1\,\% superior con probabilidad $0.22$: hay un sesgo fuerte
hacia los peores sin llegar a ser determinista.

\paragraph{Optimización: lista ordenada incremental.} La formulación directa recalcula y
ordena los costos de todos los clientes antes de cada extracción, con costo
$\bigO(q\,n\log n)$. Pero retirar un cliente solo cambia el costo de retiro de sus dos vecinos
en la ruta. El Algoritmo~\ref{alg:worst} ordena una sola vez y después mantiene la lista:
el cliente extraído y sus vecinos se localizan por búsqueda binaria, usando el orden total
(costo decreciente, identificador creciente), y se reubican. Un mapa auxiliar cliente
$\to$ ruta evita buscar al cliente en toda la solución.

\begin{algorithm}[H]
\caption{\fn{ExtracciónDePeores} (\cod{worstRemoval})}
\label{alg:worst}
\Entrada{solución $x$, grado $q$}
construir $\textit{rutaDe}[\cdot]$ (cliente $\to$ índice de ruta)\;
calcular $\operatorname{cr}(i)$ para todo cliente en ruta según~\eqref{eq:cr}\;
$\mathcal{L}\gets$ clientes ordenados por ($\operatorname{cr}$ decreciente, identificador creciente)\;
\Mientras{se hayan extraído menos de $q$ clientes \Y $\mathcal{L}\neq\emptyset$}{
  $y\sim\unif{0}{1}$;\quad $\iota\gets\min\bigl(\lfloor y^{p}|\mathcal{L}|\rfloor,\ |\mathcal{L}|-1\bigr)$\;
  \uSi{$\mathcal{L}[\iota]$ empata en costo con $\mathcal{L}[\iota-1]$ \O con $\mathcal{L}[\iota+1]$}{
    $i\gets$ \fn{ReferenciaPeores}($x$, $\iota$) \Com*[r]{ordenamiento completo original}
  }
  \lSiNo{$i\gets\mathcal{L}[\iota]$}
  $\pi\gets$ ruta $\textit{rutaDe}[i]$;\quad $k\gets$ posición de $i$ en $\pi$\;
  quitar $i$ de $\mathcal{L}$ por búsqueda binaria;\quad $\textit{rutaDe}[i]\gets -1$\;
  mover $i$ a $\Uns$;\quad \fn{RecalcularRuta}($\pi$)\;
  \ParaCada{vecino $v\in\{\pi_{k-1},\pi_k\}$ con $v\neq 0$}{
    $c'\gets\operatorname{cr}(v)$ en la ruta modificada\;
    \lSi{$c'\neq\operatorname{cr}(v)$ anterior}{reubicar $v$ en $\mathcal{L}$ por búsqueda binaria con costo $c'$}
  }
}
\end{algorithm}

\paragraph{Equivalencia con la versión de referencia.} Cuando el elemento elegido empata en
costo con alguno de sus vecinos en la lista, el orden entre ellos depende del algoritmo de
ordenamiento; en ese caso el operador delega la elección en el procedimiento original
(\cod{worstByFullSort}), que reconstruye y ordena la lista completa. Cuando no hay empate, el
elemento de índice $\iota$ es único y coincide en ambas versiones. De este modo la versión
optimizada elige siempre el mismo cliente que la original y, para una misma semilla, la
trayectoria de búsqueda es idéntica: la optimización reduce el tiempo de cómputo sin alterar
los resultados. El mismo principio se aplica en la extracción de Shaw y en los operadores de
reparación.

El costo pasa de $\bigO(q\,n\log n)$ a $\bigO(n\log n)$ para el ordenamiento inicial más, por
extracción, tres búsquedas binarias, el desplazamiento contiguo de elementos de la lista y un
recálculo de ruta $\bigO(\bar{L})$.

\subsection{Extracción de Shaw}
\label{sec:shaw}

Propuesta por Shaw~\cite{shaw1998}, extrae clientes \emph{parecidos} entre sí, con la idea de
que son intercambiables y la reparación podrá recolocarlos de otra manera. La afinidad entre
dos clientes combina distancia, diferencia de apertura de ventana y diferencia de demanda
(menor valor significa más afinidad):
\begin{equation}
\Rel(i,j)=9\,d_{ij}+3\,|a_i-a_j|+2\,|\mu_i-\mu_j| .
\label{eq:rel}
\end{equation}
Los pesos $(9,3,2)$ son los de~\cite{ropke2006}; los términos no se normalizan, por lo que su
influencia relativa depende de las escalas de la instancia. El primer cliente se elige al azar
(ruta y posición). Después, en cada paso, se toma como base un cliente al azar entre los ya
extraídos y se extrae uno de los que siguen en ruta, elegido entre los más afines a la base
con el mismo sesgo~\eqref{eq:sesgo} ($p=3$).

\paragraph{Optimización: orden de afinidad precalculado.} $\Rel(i,j)$ depende solo de la
instancia, no de la solución. Por ello el orden de afinidad de cada cliente respecto de todos
los demás puede calcularse una única vez. La primera invocación del operador construye la
tabla $\textit{orden}[b]$ (\cod{Instance::shaw\_order}), que contiene para cada cliente $b$
los otros $n-1$ clientes ordenados por (afinidad creciente, identificador creciente). El costo
es $\bigO(n^2\log n)$ una sola vez por ejecución y la memoria $n(n-1)$ enteros (unos $2.6$~MB
para $n=800$); la tabla pertenece a la instancia y la comparten todas las soluciones.

A partir de ahí no se ordena nada: el Algoritmo~\ref{alg:shaw} recorre la fila de la base
saltando los clientes que ya no están en ruta hasta alcanzar el rango deseado. Como el sesgo
concentra $\iota$ en valores pequeños, el recorrido suele ser corto.

\begin{algorithm}[H]
\caption{\fn{ExtracciónDeShaw} (\cod{shawRemoval})}
\label{alg:shaw}
\Entrada{solución $x$, grado $q$}
elegir al azar una ruta activa y un cliente $i_0$ de ella; mover $i_0$ a $\Uns$\;
$\mathcal{E}\gets\{i_0\}$ \Com*[r]{clientes extraídos por este operador}
construir $\textit{rutaDe}[\cdot]$;\quad $\textit{enRuta}\gets$ número de clientes en ruta\;
\Mientras{$|\mathcal{E}|<q$ \Y $\textit{enRuta}>0$}{
  $b\gets$ elemento al azar de $\mathcal{E}$\;
  $y\sim\unif{0}{1}$;\quad $\iota\gets\min\bigl(\lfloor y^{p}\,\textit{enRuta}\rfloor,\ \textit{enRuta}-1\bigr)$\;
  recorrer $\textit{orden}[b]$ saltando los clientes con $\textit{rutaDe}=-1$ hasta obtener los de rango $\iota-1$, $\iota$ e $\iota+1$\;
  $i\gets$ cliente de rango $\iota$\;
  \Si{$\Rel(b,i)$ empata con la del rango $\iota-1$ \O la del rango $\iota+1$}{
    $i\gets$ \fn{ReferenciaShaw}($x$, $b$, $\iota$) \Com*[r]{ordenamiento completo original}
  }
  mover $i$ a $\Uns$ y a $\mathcal{E}$;\quad \fn{RecalcularRuta}(ruta de $i$)\;
  $\textit{rutaDe}[i]\gets -1$;\quad $\textit{enRuta}\gets\textit{enRuta}-1$\;
}
\end{algorithm}

El costo por extracción pasa de $\bigO(n\log n)$ (calcular y ordenar las afinidades de todos
los clientes en ruta) a un recorrido lineal en el peor caso y típicamente mucho más corto. El
tratamiento de empates es el mismo que en la Sección~\ref{sec:worst} y garantiza la
equivalencia con la versión de referencia.

\subsection{Extracción de grupos}
\label{sec:cluster}

Es la extracción de \emph{clusters} de Pisinger y Ropke~\cite{pisinger2007}. Parte de la
observación de que, al extraer clientes parecidos de una ruta, la reparación tiende a
devolverlos a ella si el resto de la ruta sigue cerca. El operador divide los clientes de una
ruta en dos grupos geográficos y extrae uno \emph{entero}, y repite el proceso en rutas
vecinas hasta alcanzar $q$ (Algoritmo~\ref{alg:cluster}).

Los dos grupos se obtienen con el algoritmo de Kruskal sobre el grafo completo de los clientes
de la ruta, con las distancias $d_{ij}$ como pesos, \emph{detenido cuando quedan dos
componentes conexas}: las aristas se recorren de menor a mayor y cada una fusiona las
componentes de sus extremos si son distintas (estructura de conjuntos disjuntos con
compresión de caminos). El resultado equivale a quitar la arista más larga del árbol de
expansión mínima, es decir, a separar la ruta por su mayor salto espacial.

\begin{algorithm}[H]
\caption{\fn{ExtracciónDeGrupos} (\cod{clusterRemoval})}
\label{alg:cluster}
\Entrada{solución $x$, grado $q$}
$\pi\gets$ ruta activa elegida uniformemente al azar;\quad $\mathcal{E}\gets\emptyset$ \Com*[r]{extraídos por este operador}
\Mientras{exista $\pi$}{
  $\mathcal{C}\gets$ clientes de $\pi$\;
  ordenar los pares $\{a,b\}\subseteq\mathcal{C}$ por $d_{ab}$ creciente\;
  fusionar componentes en ese orden hasta que queden $2$ \Com*[r]{Kruskal truncado}
  $\mathcal{G}\gets$ componente de un cliente de $\mathcal{C}$ elegido uniformemente al azar\;
  mover los clientes de $\mathcal{G}$ a $\Uns$ y a $\mathcal{E}$;\quad \fn{RecalcularRuta}($\pi$)\;
  \lSi{$|\mathcal{E}|\ge q$}{\Salir}
  $s\gets$ elemento al azar de $\mathcal{E}$\;
  $\pi\gets$ ruta, distinta de la actual, que contiene al cliente en ruta más cercano a $s$ (ninguna si no quedan clientes en otras rutas)\;
}
\end{algorithm}

Conviene destacar tres detalles:
\begin{itemize}
\item \textbf{Elección del grupo.} El grupo se elige sorteando un cliente de la ruta, no una
de las dos componentes, por lo que la probabilidad de extraer un grupo es proporcional a su
tamaño. Si la ruta tiene uno o dos clientes, cada componente es un único cliente; con un solo
cliente, la ruta queda vacía.
\item \textbf{Encadenamiento.} La siguiente ruta es la del cliente más cercano a uno de los ya
extraídos, lo que concentra la destrucción en una región del plano: los grupos extraídos son
vecinos entre sí y pueden intercambiarse entre rutas adyacentes. Solo se excluye la ruta
recién tratada, de modo que el proceso puede volver más adelante a una ruta ya dividida y
dividir lo que quedó de ella.
\item \textbf{Grado efectivo.} El operador extrae grupos completos, así que puede superar $q$
en hasta un grupo.
\end{itemize}
Frente a la extracción de Shaw, que elige cliente a cliente por afinidad con sesgo aleatorio,
aquí la unidad extraída es un bloque espacial compacto y contiguo de una ruta. El costo por
ruta tratada es $\bigO(\bar{L}^{2}\log\bar{L})$ para ordenar los pares, más $\bigO(n)$ para
localizar la ruta siguiente.

\subsection{Extracción por ventana de tiempo}

\begin{algorithm}[H]
\caption{\fn{ExtracciónPorVentana} (\cod{timeWindowRemoval})}
\label{alg:tw}
\Entrada{solución $x$, grado $q$}
$\mathcal{C}\gets$ clientes en ruta ordenados por anchura de ventana $b_i-a_i$ creciente\;
$\mathcal{D}\gets$ los $\min(q,|\mathcal{C}|)$ primeros de $\mathcal{C}$\;
ordenar $\mathcal{D}$ por (ruta, posición) decreciente \Com*[r]{mantiene válidos los índices al borrar}
\ParaCada{$i\in\mathcal{D}$}{
  mover $i$ a $\Uns$;\quad \fn{RecalcularRuta}(ruta de $i$)\;
}
\end{algorithm}

Es un operador específico del VRPTW. Los clientes con ventanas estrechas son los más difíciles
de ubicar y los que condicionan la estructura temporal de las rutas; al extraerlos todos a la
vez, la reparación puede redecidir su asignación de forma conjunta en lugar de heredar su
posición. Es el único operador de destrucción determinista: para una solución y un $q$ dados
extrae siempre el mismo conjunto, de modo que su variabilidad proviene de $q$, del operador de
reparación con que se empareje y del ruido de este. Su costo es $\bigO(n\log n)$.
